% Source: https://artofproblemsolving.com/community/c3h3313730_some_number_theory_questions
% Stable identity: aops-topic-3313730-q02
The number of pairs is nonincreasing as $n$ increases, while $n$ strictly increases, so there can be at most one solution. For $n=10$, only $x=3,4,5,6$ contribute, giving respectively 1,2,3,4 possible values of $y$. Thus the count is 10, proving $\boxed{n=10}$ is the unique solution.
\begin{asy}
usepackage("amsmath");
size(7cm);
draw((0,0)--(7,0),Arrow);draw((0,0)--(0,8),Arrow);for(int x=0;x<=6;++x)for(int y=0;y<=7;++y){dot((x,y),x+y>=10?red:black);}draw((3,7)--(6,4),blue);label("$x+y=10$",(4.5,5.5),NE);label("$x$",(7,0),S);label("$y$",(0,8),W);label("$6$",(6,0),S);label("$7$",(0,7),W);
\end{asy}
